\scanpage{10}
\uline{Lemme 5} Supposons que pour tte \deleted{\unsure{propre}} extension finie \inserted{séparable} $K'$ de $K$, on ait $\operatorname{Br}(K')=0$. Alors on a $H^i(k,\unsure{\mathscr{K}|\mathcal{K}})=0$ pour $i\geq2$, $H$ groupe commutatif fini sur $k$. [N.B. en fait, inutile que $K$ soit fini, comme on voit sans difficulté par dévissage, cf. \deleted{(\unsure{en} \unsure{raison} des \illegible{2} \unsure{groupes} \unsure{lisses} sur $k$)}
% slanted
\margin{Ce raisonnement \unsure{ci-haut} prouve aussi que $H^i(k,U)=0$ pour $i\geq2$ si $U$ unipotent, $k$ corps quelconque. Donc pour $i\unsure{\geq|>}\unsure{2|?}$, \unsure{la seule} cohomologie qui compte est la cohomologie \unsure{étale} \unsure{des|pour} \unsure{groupes} \unsure{algébriques} \unsure{comm.}.}

$\mathscr{K}$ se dévisse en groupes qui sont, soit de t. m. sur $k$, \deleted{soit \unsure{donc} \unsure{à} \unsure{groupes} \unsure{étales}} soit isomorphes à $\alpha_p$. On sait que dans le premier cas, on a une suite exacte
\diagram{\begin{tikzcd}[ampersand replacement=\&]
0 \arrow[r,dash] \& K \arrow[r,dash] \& \unsure{G|C} \arrow[r,dash] \& G' \arrow[r,dash] \& 0
\end{tikzcd}}
avec $G,G'$ des tores, et par Tate--Nakayama on a
\[
H^i(k,G)=H^i(k,G')=0\qquad\text{pour }i>0
\]
[N.B. compte tenu que la cohomologie plate $=$ cohomologie galoisienne]. On en conclut par la suite exacte $H^i(k,K)=0$ pour $i\geq2$. Dans le cas de $\alpha_p$ il est bien connu en tous cas que $H^i(k,\alpha_p)=0$ tt $i$.

\uline{Théorème} Soit $\unsure{K|k}$ un corps tel que $\operatorname{cd}(k)\leq1$. Alors pour tt \deleted{suite exacte $e-K-G-G'$} \inserted{épim. $G-G'$ de} groupes alg sur $k$, $H^1(k,G)\longrightarrow H^1(k,G')$ est un \uline{épimorphisme}
% archive sheet number: 9

